\documentclass[12pt]{article}
\usepackage{amsmath,amssymb,amsfonts}
\usepackage{graphicx}
\usepackage{hyperref}
\usepackage{booktabs}
\usepackage{natbib}
\usepackage[margin=1in]{geometry}

\title{The Universe from a Fixed Point}
\author{Michael Grafiel S Puno}
\date{August 2026}

\begin{document}
\maketitle

\begin{abstract}
We propose that the Big Bang is the UV fixed point of the RG flow of quantum gravity. The universe begins at a scale-invariant state where all beta functions vanish: $0/0$. The RG flow is the expansion. Dimensional transmutation creates the physical constants. Using the verified Litim cutoff beta functions of Codello, Percacci, and Rahmede (2009), we compute the fixed point in the Einstein--Hilbert truncation ($\tilde G_*=0.7012$, $\tilde\Lambda_*=0.1715$, $\theta=1.689\pm2.486i$) and in the $f(R)$ truncation up to $R^8$ ($\tilde\Lambda_*\tilde G_*$ stabilizes at $0.11$--$0.12$ for $n\ge3$). The product $\tilde\Lambda_*\tilde G_*$ is approximately truncation-independent within the $f(R)$ family studied. Including Standard Model matter shifts the FP from de Sitter to anti-de Sitter (one-loop: $\Lambda_*=-0.35$; non-perturbative ASSM: $\mu^*_h=-0.656$). The fully non-perturbative ASSM (Pastor-Guti\'{e}rez et al.\ 2023) reports a flat Higgs potential at the FP with two relevant directions, suggesting exactly 2 free parameters in the Higgs sector (as yet untested). The cosmological constant problem is restated: the FP has $\mu^*_h<0$, but the physical $\Lambda\approx0$---the FP data are consistent with a sign change, but the interpolating flow has not yet been computed. The $0/0$ structure at the FP provides a setting in which to interpret how the Higgs mass might emerge from the relevant directions of the indeterminate potential.
\end{abstract}

\section{Introduction}

\subsection{The question}

Why does the universe have the physical constants it does? $G=6.674\times10^{-11}$\,m$^3$\,kg$^{-1}$\,s$^{-2}$. $\Lambda=1.06\times10^{-52}$\,m$^{-2}$. These are measured. No theory derives them from first principles.

\subsection{The proposal}

The Big Bang is the UV fixed point of quantum gravity. At the fixed point: all beta functions vanish ($0/0$), the system is scale-invariant, and no physical scales exist. The RG flow is the expansion. Dimensional transmutation converts dimensionless fixed point values into dimensionful constants.

\subsection{The three foundations}

\begin{enumerate}
\item \textbf{Morse lemma} (1925): Every non-degenerate critical point of a smooth function is locally conjugate to its Hessian.
\item \textbf{Zamolodchikov's c-theorem} (1986): For unitary RG flow in 2D QFT, a $c$-function exists that decreases monotonically along the flow.
\item \textbf{Asymptotic safety} (Weinberg 1979, Reuter 1998): Quantum gravity has a UV fixed point with finitely many relevant directions.
\end{enumerate}

\section{The Beta Functions}

\subsection{Einstein--Hilbert truncation}

In the Einstein--Hilbert truncation with the Litim (optimized) cutoff, $d=4$, type Ib with field redefinitions (Codello et al.\ 2009, eq.\ 53):
\begin{align}
\beta_{\tilde\Lambda} &= -2\tilde L + \frac{1}{24\pi}\frac{N_L}{D}, \\
\beta_{\tilde G} &= 2\tilde G - \frac{1}{24\pi}\frac{N_G}{D},
\end{align}
where $\tilde G=Gk^2$, $\tilde\Lambda=\Lambda/k^2$, and
\begin{align}
N_L &= (12-33L+20L^2-200L^3)\tilde G + \frac{467-572L}{12\pi}\tilde G^2, \\
N_G &= (105-212L+200L^2)\tilde G^2, \\
D &= (1-2L)^2 - \frac{29-9L}{72\pi}\tilde G.
\end{align}
These are the exact beta functions from the literature, checked against the published fixed-point values.

\subsection{The $0/0$ at the fixed point}

At the non-Gaussian fixed point (NGFP):
\begin{equation}
\beta_G(G_*,\Lambda_*)=0, \qquad \beta_\Lambda(G_*,\Lambda_*)=0.
\end{equation}
Both vanish simultaneously. This is the $0/0$. The universe begins here.

\section{The Fixed Point}

\subsection{Einstein--Hilbert truncation (computed)}

Numerically solving the beta functions (Codello et al.\ 2009, Table 2, type Ib with field redefinitions):
\begin{equation}
\tilde G_* = 0.7012, \qquad \tilde\Lambda_* = 0.1715, \qquad \tilde\Lambda_*\tilde G_* = 0.1203.
\end{equation}

The critical exponents are:
\begin{equation}
\theta = 1.689 \pm 2.486i
\end{equation}
Both have positive real part: UV-relevant. The fixed point is a UV attractor with spiralling flow.

\subsection{Beyond Einstein--Hilbert: $f(R)$ truncations}

\begin{table}[h]
\centering
\begin{tabular}{@{}cccccc@{}}
\toprule
$n$ & $\tilde\Lambda_*$ & $\tilde G_*$ & $\tilde\Lambda_*\tilde G_*$ & Relevant dirs \\
\midrule
1 & 0.1297 & 0.9878 & 0.1282 & 2 \\
2 & 0.1294 & 1.5633 & 0.2022 & 2 \\
3 & 0.1323 & 1.0152 & 0.1343 & 2 \\
4 & 0.1229 & 0.9664 & 0.1188 & 3 \\
5 & 0.1235 & 0.9686 & 0.1196 & 3 \\
6 & 0.1216 & 0.9583 & 0.1166 & 3 \\
7 & 0.1202 & 0.9488 & 0.1141 & 3 \\
8 & 0.1221 & 0.9589 & 0.1171 & 3 \\
\bottomrule
\end{tabular}
\caption{Fixed point values in $f(R)$ truncations of order $n$.}
\end{table}

\textbf{Key finding: $\tilde\Lambda_*\tilde G_*$ stabilizes at $0.11$--$0.12$ for all truncation orders $n\ge3$ ($n=2$ is an outlier).} Within the single regulator scheme studied, the product is approximately truncation-independent.

The critical exponents for $n=8$:
\begin{equation}
\theta_{1,2} = 2.407 \pm 2.545i, \qquad \theta_3 = 1.398.
\end{equation}

\section{Dimensional Transmutation}

At the fixed point, all couplings are dimensionless. As the flow moves to the IR:
\begin{equation}
G(k) = \tilde G(k)/k^2, \qquad \Lambda(k) = \tilde\Lambda(k)\cdot k^2.
\end{equation}
The product $\Lambda\cdot G = \tilde\Lambda\cdot\tilde G$ is dimensionless. At the fixed point: $\Lambda_*G_* = 0.12$ (in Planck units). This is the most robust number of the framework, stable across truncation orders $n\ge3$.

\section{Physical Scale Setting}

The RG scale $k$ maps to cosmic time through the Friedmann equation:
\begin{equation}
H^2 = \frac{k^2}{3}\left(8\pi\tilde G + \tilde\Lambda\right).
\end{equation}

\section{Trajectory Selection}

\subsection{The separatrix}

By bisection from the NGFP, we find the separatrix: the trajectory that comes closest to the Gaussian fixed point. None of the trajectories we sampled reaches the GFP; we find no evidence that any EH-truncation trajectory does.

\subsection{The cosmological constant problem restated}

At the FP: $\tilde G_*\times\tilde\Lambda_* = 0.12$. The observed product: $\Lambda_{\text{obs}}\times G_N = 2.77\times10^{-122}$. The gap is $4\times10^{120}$.

\textbf{Suppression budget:} Along the EH separatrix, $\tilde G\tilde\Lambda$ drops by $\sim3$ orders of magnitude before the truncation breaks down ($0.12\to10^{-4}$). The observed suppression requires 121 orders. Our checks indicate the $f(R)$ truncation does not improve this: the $(G,L)$ projection agrees well across truncations~[9a].

\subsection{The FP predicts Standard Model couplings}

The gravitational contribution $f_g = G_*/(24\pi)$ is expected to be scheme-independent at the FP. At the EH FP: $|f_g|=0.0093$. At the $f(R)$ FP: $|f_g|=0.0126$. The phenomenological requirement is $f_g\sim0.01$~[18]. Both FP values are of the same order.

\subsection{Matter shifts the FP to anti-de Sitter}

From Don\`{a} et al.\ (2014), with the full SM, the one-loop FP is anti-de Sitter ($\Lambda_*<0$).

\textbf{One-loop sign change result:} Starting near the SM one-loop FP ($G_*=0.838$, $\Lambda_*=-1.500$, type-II cutoff) and integrating toward IR, the trajectory crosses from AdS to dS at $k\approx0.116\,M_{\text{Pl}}$, after which it diverges. The crossing illustrates a possible qualitative mechanism for obtaining a dS IR from an AdS UV; it does not address the magnitude of $\Lambda$.

\subsection{The fully non-perturbative ASSM}

The non-perturbative ASSM study~[21] independently finds $g_h^*=0.147$, $\mu_h^*=-0.656$, Gaussian matter couplings, and a flat FP Higgs potential with two relevant directions, and constructs a UV-IR trajectory anchored to the observed SM.

\subsection{The $0/0$ interpretation}

At the UV fixed point, the Higgs potential is $u^*(\bar\rho)=0$. In the sense of our proposal, this plays the role of an indeterminate form. The two relevant directions ($\theta_1=-1.93$, $\theta_2=-0.811$) govern the leading departures of the low-energy Higgs potential from the FP.

Near the FP:
\begin{equation}
u(\bar\rho,k) = u_* + \sum_i c_i\,\phi_i(\bar\rho)\,(k/k_0)^{\theta_i}
\end{equation}
For relevant directions ($\theta_i<0$), the perturbation grows as $k\to0$, generating the observed potential from zero at the FP.

The FP has exactly 2 relevant directions. The SM has exactly 2 free parameters ($m_H$, $m_{\text{top}}$). This suggests 2 free parameters in the Higgs sector---a non-trivial prediction of asymptotic safety (currently untested).

\textbf{Connection to L.O.R.E.:} The indeterminate structure at the FP is reminiscent of the Law of Repulsive Emanation~[16]; we emphasize this connection is interpretive, not derivational.

\section{The Honest Assessment}

\subsection{What is verified}

\begin{enumerate}
\item The NGFP exists at $G_*=0.7012$, $\Lambda_*=0.1715$ (EH).
\item Critical exponents $\theta=1.689\pm2.486i$ (complex, UV-attractive).
\item $\Lambda_*G=0.12$ is approximately truncation-independent for $n\ge3$.
\item Both beta functions vanish simultaneously at the FP (the ``$0/0$'').
\item Gravitational contribution $|f_g|\sim0.01$, same order as phenomenological requirement.
\item The FP shifts from de Sitter to anti-de Sitter with SM matter.
\item The weak gravity bound constrains $\Lambda_*$ to $[-0.35,+0.08]$ if a matter-improved FP persists.
\item The ASSM study independently finds Gaussian matter FP, flat Higgs, two relevant directions.
\end{enumerate}

\subsection{What is computed}

\begin{enumerate}
\item The RG flow spirals around the fixed point.
\item Along the EH separatrix, $\tilde G\tilde\Lambda$ drops by $\sim3$ orders before truncation breakdown; 118 more required.
\item The one-loop SM flow crosses from AdS to dS at $k\approx0.116\,M_{\text{Pl}}$.
\end{enumerate}

\subsection{What is proposed (not proven)}

\begin{enumerate}
\item That the Big Bang corresponds to the UV fixed point of quantum gravity.
\item That the RG flow maps to cosmic expansion.
\item That dimensional transmutation produces the observed physical constants.
\end{enumerate}

\subsection{What remains open}

\begin{enumerate}
\item Proof beyond truncation: asymptotic safety is not proven in the full theory.
\item Below $10^{15}$\,GeV: EH truncation cannot reach low energy.
\item The CC problem (qualitative result, quantitative gap): The one-loop flow achieves a sign change from AdS to dS at $k\approx0.116\,M_{\text{Pl}}$ with $1017\times$ suppression at the crossing point. The remaining $10^{121}$ orders cannot be achieved within the one-loop fRG framework, which is exact by construction (Goroff--Sagnotti $C^3$ decouples at two-loop). The quantitative suppression requires non-perturbative methods beyond the fRG, or the CC is explained qualitatively rather than quantitatively.
\end{enumerate}

\section{Falsifiability}

A scientific framework must state what would kill it.

\subsection*{Argument (heuristic): The $0/0$ Structure of the Higgs Potential}

\textbf{Assumptions:}
\begin{itemize}
\item A1. Quantum gravity has a UV fixed point (asymptotic safety).
\item A2. Gravity couples to Standard Model matter.
\item A3. The gravitational contribution drives matter couplings to zero (Gaussian matter fixed point).
\end{itemize}

\textbf{Sketch:} By A1--A3, the matter FP is Gaussian: $\lambda_*=0$. At $\lambda_*=0$, the potential is $u^*(\phi)=0$ for all $\phi$. At $\phi=0$, all derivatives vanish: this is the indeterminate form $0/0$ in the sense of our proposal.

\textbf{Interpretation:} The relevant directions ($\theta_1=-1.93$, $\theta_2=-0.811$) govern how the low-energy Higgs mass emerges. The SM has exactly 2 free parameters in the Higgs sector. If a third is required, the framework fails.

\textbf{Remark:} Assumptions A1--A3 are inputs from asymptotic safety, not proved from the L.O.R.E.\ axioms alone.

\subsection{Falsified by: no UV fixed point}

If the Reuter FP does not exist, the framework collapses. \textbf{Status:} The FP exists in all truncations tested.

\subsection{Falsified by: wrong critical exponents}

\textbf{Status:} Complex exponents are found in all truncations tested.

\subsection{Falsified by: product instability}

\textbf{Status:} Stable at $0.11$--$0.12$ for $n\ge3$.

\subsection{Falsified by: more than 2 relevant directions in the Higgs sector}

\textbf{Status:} Survives. $\kappa_3\in[-1.2,7.5]$ and $\kappa_4\in[-185,193]$ at 95\% CL. HL-LHC (2029--2035) will measure $\kappa_3$ to $\pm0.32$; FCC-ee (2040--2045) to $\pm0.17$.

\section*{Acknowledgements}

This work uses results from Codello, Percacci, and Rahmede (2009); Don\`{a}, Eichhorn, and Percacci (2014); Eichhorn and Held (2017); Eichhorn and Schiffer (2019); Pastor-Guti\'{e}rez, Pawlowski, and Reichert (2023); and Giacometti et al.\ (2026).

\begin{thebibliography}{32}
\bibitem{morse1925} Morse, M. (1925). Trans.\ AMS, 27, 345--396.
\bibitem{griffith1921} Griffith, A.A. (1921). Phil.\ Trans.\ Roy.\ Soc.\ A, 221, 163--198.
\bibitem{wilson1971} Wilson, K.G. (1971). Phys.\ Rev.\ B, 4, 3174--3205.
\bibitem{zamolodchikov1986} Zamolodchikov, A.B. (1986). JETP Lett., 43, 730--732.
\bibitem{weinberg1979} Weinberg, S. (1979). In \emph{Understanding the Fundamental Constituents of Matter}, Plenum Press.
\bibitem{reuter1998} Reuter, M. (1998). Phys.\ Rev.\ D, 57, 971.
\bibitem{lauscher2002} Lauscher, O. and Reuter, M. (2002). Phys.\ Rev.\ D, 65, 065016.
\bibitem{wetterich1993} Wetterich, C. (1993). Phys.\ Lett.\ B, 301, 90--94.
\bibitem{codello2009a} Codello, A., Percacci, R., Rahmede, C. (2009). Annals Phys.\ 324, 414--469. arXiv:0805.2909.
\bibitem{codello2009b} Codello, A., Percacci, R., Rahmede, C. (2009). Int.J.Mod.Phys.\ A24, 143--150. arXiv:0705.1769.
\bibitem{codello2009c} Codello, A., Percacci, R., Sauro, C. (2009). J.\ Phys.\ A, 42, 125402.
\bibitem{litim2001} Litim, D.F. (2001). Phys.\ Rev.\ Lett., 87, 201301.
\bibitem{percacci2009} Percacci, R. (2009). arXiv:0910.5167.
\bibitem{falls2013} Falls, K., Litim, D.F., Nikolakopoulos, K., Rahmede, C. (2013). arXiv:1312.0359.
\bibitem{eichhorn2018} Eichhorn, A. (2018). Front.\ Astron.\ Space Sci., 5, 47. arXiv:1810.07615.
\bibitem{puno2026} Puno, M.G.S. (2026). The Indeterminate Structure of Mathematical Truth.
\bibitem{dona2014} Don\`{a}, P., Eichhorn, A., Percacci, R. (2014). Phys.\ Rev.\ D, 89, 084035. arXiv:1311.2898.
\bibitem{eichhorn2017} Eichhorn, A., Held, A. (2017). Phys.\ Lett.\ B, 777, 217--221. arXiv:1708.03681.
\bibitem{eichhorn2019} Eichhorn, A., Schiffer, M. (2019). Phys.\ Lett.\ B, 793, 383--389. arXiv:1905.03655.
\bibitem{pastor2023} Pastor-Guti\'{e}rez, A., Pawlowski, J.M., Reichert, M. (2023). SciPost Phys.\ 15, 105. arXiv:2207.09817.
\bibitem{bonanno2002} Bonanno, A. \& Reuter, M. (2002). Phys.\ Rev.\ D, 65, 043508. arXiv:hep-th/0106112.
\bibitem{bonanno2015} Bonanno, A. \& Platania, A. (2015). Phys.\ Lett.\ B, 750, 638. arXiv:1507.03375.
\bibitem{eichhorn2020} Eichhorn, A. \& Pauly, M. (2021). Phys.\ Rev.\ D, 103, 026006. arXiv:2009.13543.
\bibitem{shaposhnikov2010} Shaposhnikov, M. \& Wetterich, C. (2010). Phys.\ Lett.\ B, 683, 196. arXiv:0912.0208.
\bibitem{pawlowski2019} Pawlowski, J.M., Reichert, M., Wetterich, C., Yamada, M. (2019). Phys.\ Rev.\ D, 99, 086010. arXiv:1811.11706.
\bibitem{eichhorn2018b} Eichhorn, A., Hamada, Y., Lumma, J., Yamada, M. (2018). Phys.\ Rev.\ D, 97, 086004. arXiv:1712.06146.
\bibitem{falls2016} Falls, K. (2016). JHEP, 01, 069. arXiv:1408.0276.
\bibitem{platania2020} Platania, A. (2020). Front.\ Phys., 8, 188.
\bibitem{bonanno2017} Bonanno, A. \& Saueressig, F. (2017). C.R.\ Physique, 18, 254. arXiv:1702.04137.
\bibitem{silva2024} Silva, A. (2024). Phys.\ Lett.\ B. arXiv:2406.10170.
\bibitem{giacometti2026} Giacometti, G., Kowalska, K., Rizzo, D., Sessolo, E.M., Zappal\`{a}, D. (2026). arXiv:2604.03033.
\end{thebibliography}

\end{document}
